% ============================================================
% 附录 A：后端 REST 接口一览
% ============================================================
\appendixsection{后端 RESTful 接口一览}\label{app:api}

Web Console 后端采用 RESTful 设计，完整接口列表如下：

\begin{table}[htbp]
\centering
  \renewcommand{\arraystretch}{1.3}
\begin{tabular}{lll}
    \toprule
\textbf{方法} & \textbf{路径} & \textbf{功能} \\
    \midrule
GET  & /api/config                     & 读取系统配置 \\
    \midrule
POST & /api/config                     & 更新系统配置 \\
    \midrule
GET  & /api/vulns                      & 获取漏洞列表 \\
    \midrule
POST & /api/task/start                 & 创建新的渗透测试任务 \\
    \midrule
GET  & /api/task/status                & 获取当前任务状态 \\
    \midrule
POST & /api/task/stop                  & 停止当前任务 \\
    \midrule
GET  & /api/task/logs/<id>?offset=N    & 增量拉取指定任务的执行日志 \\
    \midrule
GET  & /api/task/history               & 获取任务历史列表 \\
    \midrule
GET  & /api/docker/status              & 获取 Docker 运行状态与容器列表 \\
    \midrule
GET  & /api/docker/envs                & 列出可用的 Docker Compose 环境 \\
    \midrule
POST & /api/docker/start               & 启动 Docker Compose 环境 \\
    \midrule
POST & /api/docker/stop                & 停止 Docker Compose 环境 \\
    \midrule
GET  & /api/results                    & 获取测试结果列表 \\
    \midrule
GET  & /api/results/stats              & 获取统计汇总数据 \\
    \midrule
GET  & /api/system/info                & 获取系统运行信息 \\
    \bottomrule
\end{tabular}
\end{table}

% ============================================================
% 附录 B：项目目录结构
% ============================================================
\appendixsection{项目目录结构}\label{app:structure}

以下为 AutoPT+ 项目的完整目录结构，来源于项目根目录的 \texttt{README.md}：

\begin{Verbatim}[mathescape=false]
AutoPT/
+-- app.py                        # Flask Web后端
+-- install_tools.ps1             # Windows扫描工具一键安装脚本
+-- frontend/                     # Web前端界面
|   +-- index.html                # 主页面（SPA）
|   +-- style.css                 # 样式（暗色主题）
|   +-- app.js                    # 前端逻辑（路由/API/交互）
+-- src/                          # 渗透测试引擎核心
|   +-- main.py                   # CLI入口
|   +-- autopt.py                 # AutoPT核心类（状态机构建/运行）
|   +-- prompt.py                 # Agent提示词模板（ReAct格式）
|   +-- react_parser.py           # 自定义ReAct输出解析器
|   +-- tools.py                  # Agent工具定义
|   +-- terminal.py               # 命令执行（本地/SSH远程）
|   +-- knowledge.py              # 服务知识库加载
|   +-- utils.py                  # 工具函数
|   +-- config/
|   |   +-- config.yml            # 项目主配置文件
|   |   +-- services.yml          # 服务知识库
|   +-- psm/                      # 渗透状态机模块
|   |   +-- __init__.py
|   |   +-- state.py              # 状态定义 + 上下文压缩 + Check验证
|   |   +-- trans.py              # 状态转换路由
|   |   +-- utils.py              # 辅助函数
|   +-- result/                   # 测试结果输出目录
+-- bench/                        # 测试基准数据集
|   +-- data.jsonl                # 20个CVE漏洞元数据
|   +-- [漏洞类型]/[CVE编号]/
|       +-- docker-compose.yml    # 各漏洞的Docker环境
+-- xray/                         # xray扫描器
|   +-- xray_windows_amd64.exe
+-- README.md
\end{Verbatim}

% ============================================================
% 附录 C：前端路由代码片段
% ============================================================
\appendixsection{前端路由代码片段}\label{app:router}

以下为 \texttt{frontend/app.js} 中实现单页应用路由切换的核心代码。系统通过 \texttt{PAGE\_TITLES} 和 \texttt{PAGE\_LOADERS} 两个映射表将页面名称与数据加载函数绑定，点击导航项时调用 \texttt{switchPage} 函数切换显示区域并触发对应的数据拉取逻辑：

\begin{Verbatim}[mathescape=false]
const PAGE_TITLES = {
    dashboard: '控制台', vulns: '漏洞库', attack: '渗透测试',
    docker: '靶机管理', results: '测试报告', settings: '系统设置'
};
const PAGE_LOADERS = {
    dashboard: loadDashboard, vulns: loadVulns,
    docker: () => { refreshDocker(); loadDockerEnvs(); },
    results: loadResults,
    settings: () => { loadConfig(); loadSystemInfo(); }
};
function switchPage(name) {
    document.querySelectorAll('.page').forEach(
        p => p.classList.remove('active'));
    document.querySelectorAll('.nav-item').forEach(
        n => n.classList.remove('active'));
    $(`page-${name}`)?.classList.add('active');
    document.querySelector(
        `.nav-item[data-page="${name}"]`)?.classList.add('active');
    $('pageTitle').textContent = PAGE_TITLES[name] || name;
    PAGE_LOADERS[name]?.();
}
document.querySelectorAll('.nav-item').forEach(
    item => item.addEventListener(
        'click', () => switchPage(item.dataset.page))
);
\end{Verbatim}

